% © 2026 Ing. David Jaroš — CC BY-NC-ND 4.0
%
% This work is licensed under a Creative Commons Attribution-NonCommercial-NoDerivatives
% 4.0 International License (CC BY-NC-ND 4.0).
%
% License History: Earlier drafts (up to v0.3) were released under CC BY 4.0.
% From v0.4 onward, all material is released under CC BY-NC-ND 4.0 to protect
% the integrity of the theoretical work during ongoing academic development.
%
% See LICENSE.md for full license text.
%
% File: canonical/alpha/alpha_gap_closure_matrix.tex
% Purpose: Single canonical record of the T3_ALPHA chain status.
%          Identifies the minimal missing theorem (Gap G137-B).
%
% Status: 2026-05-11

\ifdefined\INCLUDEMODE
\else
  \documentclass[12pt,a4paper]{article}
  \usepackage[a4paper, margin=2.5cm]{geometry}
  \usepackage{amsmath, amssymb, amsthm}
  \usepackage{mathtools}
  \usepackage{hyperref}
  \usepackage{xcolor}
  \usepackage{booktabs}
  \usepackage{longtable}
  \usepackage{tcolorbox}
  \tcbuselibrary{skins,breakable}

  \newcommand{\PROVED}{\textbf{PROVED}}
  \newcommand{\MC}{\textbf{MC}}

  \title{\textbf{Alpha Gap Closure Matrix}\\
         \large Single Canonical Record of the T3\_ALPHA Chain Status}
  \author{Ing. David Jaro\v{s}}
  \date{2026-05-11}

  \begin{document}
  \maketitle
\fi

%──────────────────────────────────────────────────────────────
\section{Purpose}
%──────────────────────────────────────────────────────────────

This document is the single canonical record of the alpha program chain
status, connecting all partial mechanisms and identifying the minimal
missing theorem.

%──────────────────────────────────────────────────────────────
\section{The Candidate Chain}
%──────────────────────────────────────────────────────────────

The full candidate derivation chain is:
\[
S[\Theta]
\;\to\; T^2\ \text{winding}
\;\to\; n\ln n
\;\to\; N_{\rm eff}
\;\to\; c{=}3
\;\to\; \eta(i)
\;\to\; B \approx 46.28
\;\to\; n^*{=}137
\;\to\; \alpha^{-1}\approx 137.
\]

\begin{center}
\begin{longtable}{p{3.5cm} p{5cm} p{2cm} p{3cm}}
\toprule
\textbf{Chain step} & \textbf{Claim} & \textbf{Status} & \textbf{Source} \\
\midrule
$S[\Theta]$ defined & UBT action with $\mathcal{L}_{kin}+\mathcal{L}_{grav}+\mathcal{L}_{gauge}$
  & \PROVED\ [L1] & \texttt{canonical\_action.tex} \\[4pt]
$T^3$ compactification & Three compact Im$(\mathbb{H})$ circles
  & \PROVED\ [L0] & \texttt{step1\_mode\_decomp.} \\[4pt]
$N_{\rm phases}=3$ & From $\dim_\mathbb{R}\,\mathrm{Im}\,\mathbb{H}$
  & \PROVED\ [L0] & \texttt{step2\_AUDIT.tex} \\[4pt]
$N_{\rm eff}^{\rm loop}=3$ & Minimal scalar loop from $S_{\rm kin}[\Theta]$
  & \PROVED\ [L1] & \texttt{neff\_reconciliation.tex} \\[4pt]
$N_{\rm eff}^{\rm loop}=12$ & KK on $T^3$: jumps $1\to7\to19$, not 12
  & \textbf{FAILED} & \texttt{neff\_kk\_counting.tex} \\[4pt]
$N_{\rm eff}=12$ (algebraic) & $\dim G_{\rm SM}=8+3+1=12$ from $\mathbb{C}\otimes\mathbb{H}$
  & \PROVED\ [L0] & \texttt{neff12\_derivations.tex} \\[4pt]
$n\ln n$ form of $V_{\rm eff}$ & Divisor sum on $T^2\subset T^3$, equal mass
  & \MC\ (math [STD]) & \texttt{nlogn\_mechanism.tex} \\[4pt]
$c=3$ CFT central charge & Mode-count from $\delta^2 S[\Theta]/\delta\Theta^2$
  & \MC & \texttt{central\_charge\_deriv.tex} \\[4pt]
$B=N_{\rm eff}^{3/2}(2\eta(i))^{c/12}$ & $=46.281$; error $0.0066\%$ vs $B_{\rm req}$
  & \textbf{OBSERVATION} & \texttt{eta\_i\_status.md} \\[4pt]
$B_0^{\rm vac.pol.}=B^{V_{\rm eff}}$ & Bridge between two distinct $B$'s
  & \textbf{OPEN} & \texttt{B\_identification\_proof.tex} \\[4pt]
$n^*(B_{\rm phenom})=137$ & From stationarity of $V_{\rm eff}$, given $B$
  & \PROVED\ [L1 given $B$] & \texttt{alpha\_best\_route.tex} \\[4pt]
$g=11$, $137\equiv1\pmod4$ unique & Direct point count $a_{131}=-9$, $a_{139}=-3$
  & \PROVED\ [L0] & \texttt{hecke\_twin\_prime\_test.md} \\[4pt]
$C\leftrightarrow w_p$ physical & Direct route failed; modular route open
  & \textbf{NO-GO/OPEN} & \texttt{charge\_conj\_wp\_bridge.tex} \\[4pt]
$f(137)\approx\Delta\alpha^{-1}_{\rm exp}$ & $0.036541$ vs $0.035999$, ratio $1.015$
  & \MC\ (not [NC]) & \texttt{alpha\_residuum\_analysis.tex} \\
\bottomrule
\end{longtable}
\end{center}

%──────────────────────────────────────────────────────────────
\section{Gap Analysis}
%──────────────────────────────────────────────────────────────

\textbf{Gap G137-B (primární)}: $B$ z vakuové polarizace a $B$ ve $V_{\rm eff}$
jsou odvozeny dvěma zcela odlišnými výpočty --- jedna z fotonového 2-pointu
(running coupling), druhá z efektivní winding akce.
Žádný teorém v kanonických souborech netvrdí že jsou si rovny.
Dokud tento most neexistuje, celý řetěz závisí na $B = B_{\rm phenom}$ jako
externím vstupu.

\medskip
\textbf{Gap G-Neff (sekundární)}: $N_{\rm eff}^{\rm loop} = 12$ nevychází
z $T^3$ KK (numericky vyloučeno). $S^1$ s $R\cdot\Lambda\in[2,3)$ dává [MC]
ale fyzikální škála není odvozena. Algebraické $N_{\rm eff}=12=\dim(G_{\rm SM})$
je jiná veličina.

\medskip
\textbf{Gap G-c3 (terciární)}: $c=3$ z mode-counting [MC] --- CFT identifikace
a RG-flow argument chybí pro povýšení na [L1].

\medskip
\textbf{Gap G-eta (kvartální)}: Odkud pochází exponent $1/4 = c/12$?
Propojení $Z_{\rm wind}=\theta_3$ je [L1] (viz \texttt{ncg\_poisson\_b0derivation.tex}),
ale identifikace konkrétní normalizační konstanty $\eta(i)$ jako fyzikálně nucené
hodnoty je [OPEN].

%──────────────────────────────────────────────────────────────
\section{The Minimal Missing Theorem}
%──────────────────────────────────────────────────────────────

\begin{tcolorbox}[colback=red!4!white, colframe=red!60!black,
  title={Minimal Missing Theorem (Gap G137-B)}]
\textbf{Statement needed}:
There exists a mechanism in $S[\Theta]$ that identifies the coefficient $B$
in the winding-mode effective potential
\[
V_{\rm eff}(n) = n^2 - B\,n\ln n
\]
with the ultraviolet coupling coefficient from the one-loop photon
self-energy, and gives
\[
B = N_{\rm eff}^{3/2}\,(2\eta(i))^{c/12} \approx 46.281
\]
from first principles, without using $\alpha_{\rm exp}$, $137$,
or $B_{\rm phenom}$ as inputs (Rules R1--R5).

\medskip
\textbf{Current status}: OPEN.
\textbf{Why this is the minimum}: all other steps in the chain are either
proved, proved given $B$, or are secondary gaps that collapse into G137-B.
\end{tcolorbox}

%──────────────────────────────────────────────────────────────
\section{What Would Constitute a Proof}
%──────────────────────────────────────────────────────────────

Three acceptable routes:

\begin{enumerate}

\item \textbf{Spectral action route}: Show that the Seeley--DeWitt coefficient
$a_4$ from the UBT Dirac operator $D$ on $M^4\times S^1$ gives precisely
$B = N_{\rm eff}^{3/2}\cdot(2\eta(i))^{1/4}$ as the normalisation of the
logarithmic term, where $N_{\rm eff}$ and the exponent $1/4$ arise from the
algebra $\mathbb{C}\otimes\mathbb{H}$.

\item \textbf{WZW route}: Establish the existence of a WZW term in $S[\Theta]$
with level $k$ derived from the gauge anomaly, and show via the Verlinde formula
that $B = \mu(\Gamma_0(p^*))/3 = (p^*+1)/3$.

\item \textbf{Partition function route}: Show that $Z_{\rm wind}(\tau)=\theta_3(0|i\tau)$
evaluated at the physical modular point $\tau^*$ (not necessarily $\tau=i$)
gives the normalisation of the effective action consistent with $B\approx 46.28$.

\end{enumerate}

%──────────────────────────────────────────────────────────────
\section{What Would Kill the Route}
%──────────────────────────────────────────────────────────────

If any of the following is established, the entire alpha route is terminated:

\begin{itemize}
  \item \textbf{K1}: Prokáže se že $B_0^{\rm vac.pol.}$ a $B^{V_{\rm eff}}$
        nemohou být sjednoceny v rámci $S[\Theta]$ (hard no-go theorem).
  \item \textbf{K2}: Prokáže se že $N_{\rm eff}=12$ nemůže vyjít z žádného
        přirozeného KK schématu na $T^3$ (zahrnuje výsledek z
        \texttt{neff\_kk\_counting.tex}: $T^3$ jump je $1\to7\to19$).
  \item \textbf{K3}: Prokáže se že $V_{\rm eff}(n)=n^2-Bn\ln n$ má jiný
        tvar pro konkrétní $S[\Theta]$ (např. $n^2-N_{\rm eff}\ln(2\sinh\pi n)$
        jako $V_{\rm CW}$ bez $n\ln n$ forma).
  \item \textbf{K4}: $|a_{137}(k=4)|=2274$ nebo $|a_{137}(k=6)|=38286$
        se prokáže jako chybná hodnota --- lepton mass signal kolabuje.
\end{itemize}

\ifdefined\INCLUDEMODE
\else
  \end{document}
\fi
